\section{Method}

\subsection{Ontology basis}

Each experiment uses a per-head ontology basis $B_{\mathrm{ont}} \in \mathbb{R}^{L \times H_{\mathrm{kv}} \times d \times r}$ constructed from the tool catalog. Intuitively, the basis spans directions that are repeatedly used by tool facets such as finance, news, travel, or media. The basis is fixed at evaluation time and shared by both the Q-side and K-side interventions.

\subsection{Practical query-space coverage bias}

For a query tensor $q$, the implemented Q-side intervention is
\[
q' = q + \beta B_{\mathrm{ont}} B_{\mathrm{ont}}^\top q.
\]
When $\beta < 0$, the operator subtracts the query component aligned with the ontology subspace. In code this is implemented as a forward hook on \texttt{q\_proj}. The method is history-free: it does not explicitly track which tool has already been emitted. We therefore interpret it as a practical approximation to a stronger ideal coverage operator, not as literal per-step facet-state tracking.

\subsection{Stationary key-side bias}

The key-side comparison uses the same basis but a different site:
\[
k' = k + \alpha B_{\mathrm{ont}} B_{\mathrm{ont}}^\top k.
\]
This intervention is stationary by construction. It amplifies the same ontology-aligned key directions at every decoding step. If the first emitted tool already occupies the dominant facet, the operator has no mechanism to suppress it at later steps.

\subsection{Working hypothesis}

The paper's working hypothesis has two parts. The first is an \emph{accuracy hypothesis}: negative query-side ontology bias should improve multi-tool recovery because it reduces repeated use of the same dominant facet. The second is a \emph{specificity hypothesis}: if the ontology basis is replaced by a random or feature-shuffled basis of the same rank, the effect should disappear. The experiments are designed to test these two claims directly.
